\section{Файловая и модульная структура}
\begin{lstlisting}
CMakeLists.txt
README.md
src/
  main.cpp
  core/
    calendarevent.h
    calendarsetup.h/.cpp
    calendarservice.h/.cpp
    simulationengine.h/.cpp
    conflictresolver.h/.cpp
    reminderservice.h/.cpp
    jsoncodec.h/.cpp
    simulationstore.h/.cpp
    simulationreport.h/.cpp
    experimentrunner.h/.cpp
    experimentmain.cpp
  ui/
    widget.h/.cpp
    simulationwindow.h/.cpp
    officescene.h/.cpp
    runtools.h/.cpp
    experimentdialog.h/.cpp
tests/
  calendarservice_test.cpp
  simulationengine_test.cpp
  completion_test.cpp
  calendarwindow_test.cpp
  fixtures/
documentation/
  README.md
  MODEL.md
  ARCHITECTURE.md
  report/
    report.tex
    sections/
    figures/
    make_results.py
    results_table.tex
    results_text.tex
    data_manifest.json
  experiments/
    results-v7/
    verification-v7/
\end{lstlisting}

\begin{longtable}{L{0.35\linewidth}L{0.57\linewidth}}
\toprule Файл или группа & Ответственность \\\midrule\endhead
\file{calendarevent.h} & Структуры данных компании, событий, фильтров, истории, контекста и результата операций. \\
\file{calendarsetup.*} & Параметры опыта, создание исходного каталога, преобразование модельной даты и времени. \\
\file{calendarservice.*} & Проверки ресурсов и прав, активный календарь, архив, история, переходы статусов. \\
\file{simulationengine.*} & Начальный календарь, модельная минута, случайные сценарии, заявки и команды. \\
\file{conflictresolver.*}, \file{reminderservice.*} & Варианты переносов, планировщик уведомлений и имитируемый канал. \\
\file{jsoncodec.*}, \file{simulationstore.*} & JSON, проверка файлов, сохранение и контролируемое воспроизведение. \\
\file{simulationreport.*} & Метрики, полный JSON, CSV и текстовая сводка. \\
\file{experimentrunner.*}, \file{experimentmain.cpp} & Серии опытов, агрегирование и консольная точка входа. \\
\file{src/ui} & Только представление и управление через публичные операции ядра. \\
\file{tests} & Четыре набора автоматических проверок. \\
\file{documentation/experiments/results-v7} & Фактические наблюдения текущего исследования, настройки, пример отчёта и сохранения. \\
\file{documentation/experiments/verification-v7} & Логи сборки и проверки, JUnit XML и сведения о среде. \\
\bottomrule
\end{longtable}

\code{calendar_core} зависит от Qt Core; \code{calendar_ui} --- от ядра и Qt Widgets. \code{Modeling} запускает GUI через \file{src/main.cpp}. \code{modeling_experiments} работает с \code{QCoreApplication} и не требует дисплея. Локальные \file{build/} и \file{.qtcreator/} содержат служебные данные и исключаются из Git. Каталог \file{documentation/experiments/results} содержит исторические измерения v4 и не подменяет набор \file{results-v7}.
